\section{Conclusion}
\label{sec:conclusion}

This paper proposes frequent losers---firms that never win yet
participate in abnormally many public tenders---as a screening marker for
cover bidding in procurement auctions. Using 4.5 million tender-items
from Brazil's BEC platform (2009--2019), we test two hypotheses of
increasing strength.

\paragraph{H1: FL as screening marker.}
We establish that FL presence is a reliable empirical marker for
procurement anomalies. FL-present tenders exhibit 4--9\% higher % CO-AUTHOR EDIT: consistent with abstract/intro
negotiated prices (OLS), robust to matching, sensitivity analysis,
alternative FL definitions, and sample restrictions. External validation
against CADE cartel convictions confirms that FL firms are 3.5 times
more likely than comparable always-losers to co-participate with known
cartelists ($p < 0.001$, permutation test), and the FL--price
association persists after excluding all CADE-involved markets (0.062
vs.\ 0.064 baseline). These findings support H1 decisively: the FL
screen identifies tenders and firms with characteristics consistent with
collusion.

\paragraph{H2: FL as cover bidders.}
We provide substantial evidence for the causal mechanism. A leave-one-out
instrumental variable strategy yields a price markup of 21\%
(first-stage $F = 396$), exceeding OLS by a factor consistent with
measurement-error attenuation in the binary treatment indicator.
Co-bidding network analysis reveals that the FL--price effect
concentrates in competitive markets---where cover bidding is most
profitable---rather than in concentrated markets where dominant firms
already control pricing. Bajari--Ye tests reject exchangeability and
conditional independence for FL bid residuals, with FL pairs exhibiting
significantly stronger within-tender correlation than non-FL pairs
(bootstrap $p < 0.001$). The evidence is consistent with Regime~1
(complementary cover bidding) where FL firms submit bids from a wide
distribution above the winning price. While we cannot conclusively prove
that every FL firm is a deliberate cover bidder, the convergence of IV
and bid-level evidence makes H2 the most parsimonious explanation.

\paragraph{Caveats and limitations.}
The staggered DiD exercise yields imprecise results, reflecting the
difficulty of identifying dynamic treatment effects in markets with
infrequent FL entry and exit. The FL definition is based on a
statistical threshold (median + 1.5 $\times$ IQR) that necessarily
involves some misclassification---though robustness to alternative
thresholds, low-win-rate variants, and cross-fitting suggests this is
not driving the results. As with any cartel screen, strategic adaptation
is a concern: if FL screening were widely adopted, cartels could rotate
cover bidders or allow them occasional wins to avoid detection.

\paragraph{Policy implications.}
The FL screen has three operational advantages for procurement agencies,
particularly in developing countries: (i) it requires only participation
and outcome data, not bid values; (ii) it produces firm-level flags,
enabling investigation of all tenders involving flagged firms; (iii) it
can be combined with bid-level screens in a two-stage workflow---FL
screening first, followed by targeted Bajari--Ye or variance-based
analysis of flagged tenders. The estimated welfare loss of R\$734 million (OLS lower bound) to
R\$2.4 billion (IV upper bound) across 2009--2019 underscores the
potential return on investment from systematic screening.

\paragraph{Future research.}
Four directions merit investigation: (i) external validation on
procurement platforms beyond BEC; (ii) deeper network analysis exploring
shared legal representatives, addresses, and financial links among FL
firms and their frequent co-winners; (iii) dynamic models of how FL
screens interact with leniency programs and cartel stability; (iv)
field experiments where procurement agencies adopt FL screening and
monitor outcomes.
